Este capítulo describe el protocolo con que se evalúa el pipeline RAG de JobAI. Los detalles de procedimiento que no se necesitan para leer los resultados están en el Anexo~\ref{anx:protocolo}.

\section{Enfoque}
\label{sec:met-enfoque}

El trabajo es un estudio de caso: el pipeline de JobAI se mejora en dos iteraciones y, al final de cada una, se evalúa \textit{offline} con una colección de prueba, de modo que el protocolo es el instrumento de medición y no el objeto evaluado. La primera iteración corrige el sistema para que haga lo que su documentación describe, y la segunda incorpora los componentes de recuperación cuyo aporte se mide (Capítulo~\ref{cap:artefacto}). Aplicar el protocolo incluye prepararlo: antes de medir se inspeccionan los planes de ejecución de las consultas, los datos y el código frente a la documentación, para que lo medido corresponda a lo que el sistema dice hacer. Toda decisión de configuración se toma con un conjunto de consultas de desarrollo, separado del de prueba, para no ajustar el sistema a las consultas con que se evalúa.

La configuración final se juzga con tres criterios: (a) que supere a BM25 en nDCG@10 con una diferencia significativa tras la corrección de Holm; (b) que cumpla las restricciones de la consulta, medido con el cumplimiento en las cinco primeras posiciones (C@5), y (c) que el retraso que agregan los componentes de la Iteración~2 hasta el primer \textit{token}, frente a la Iteración~1 y con la misma definición de latencia, quede por debajo de los 1.000~ms a partir de los cuales los usuarios de un buscador web notan un retraso con mucha probabilidad \citep{Arapakis2014,Bai2017}. Estos criterios se formularon a partir de lo que el protocolo ya medía y no son umbrales fijados antes de medir; sí se fijaron antes de la corrida final las reglas de decisión del conjunto de desarrollo, la corrección por comparaciones múltiples y la unidad de análisis.

\section{Preguntas de evaluación}
\label{sec:met-preguntas}

El cuarto objetivo específico se operacionaliza en cuatro preguntas de evaluación (PE). El motor de almacenamiento vectorial se incluye por el \enquote{con qué costo} de la pregunta de investigación, porque fija el costo de la etapa de búsqueda; su efecto en la calidad se mide en PE1. El tercer objetivo se responde con las fallas que revela la aplicación del protocolo y con el aporte de los componentes de la Iteración~2.

\begin{description}
    \item[PE1 (calidad de recuperación).] ¿La recuperación de JobAI obtiene mayor precisión y mejor ordenamiento de resultados que la búsqueda léxica, en qué tipos de consulta se concentra la diferencia y cuánto aporta cada componente del pipeline?
    \item[PE2 (fidelidad de la generación).] ¿En qué medida las respuestas que genera el pipeline se sostienen en las ofertas recuperadas y responden a la pregunta formulada, y cómo cambian entre la Iteración~1 y la 2?
    \item[PE3 (latencia).] ¿Cuánto tarda cada etapa del pipeline y cuál es la latencia percibida por el usuario?
    \item[PE4 (almacenamiento vectorial).] ¿Cómo se comparan pgvector y Qdrant en la relación entre \textit{recall} y latencia, con y sin filtros estructurados y a distintas escalas de datos?
\end{description}

\section{Corpus y consultas}
\label{sec:met-corpus}
\label{sec:met-consultas}
\label{sec:met-sistema}

El corpus se obtiene con una ejecución completa de la ingesta de JobAI sobre sus tres fuentes públicas, con la ingesta periódica desactivada, y se congela en una copia de seguridad (\textit{snapshot}) identificada por su resumen SHA-256, sobre la que se ejecutan todas las corridas. Como las fuentes están orientadas al empleo remoto pero no todas sus ofertas lo son, la modalidad se trata como una condición que la consulta puede exigir y que se evalúa, no como un rasgo de todo el corpus.

El conjunto de prueba tiene 75 consultas en inglés, el idioma de las ofertas, repartidas en las cinco categorías de la Tabla~\ref{tab:met-categorias}. Cada consulta se asocia a una \emph{necesidad de información}, una descripción breve de lo que el usuario busca. Se suman 25 traducciones al español de consultas en inglés, hechas por el autor, que reutilizan los juicios de su contraparte porque la relevancia de una oferta depende de la necesidad y no del idioma. Así, 85 consultas producen una lista de ofertas (60 en inglés y 25 en español) y 15 son de agregación. Las 24 consultas que expresan restricciones, o que podrían leerse así, se anotan con los filtros que deberían extraerse de ellas, lo que permite medir la extracción y el cumplimiento de las restricciones sin intervención del juez.

\begin{table}[htbp!]
    \centering
    \caption{Categorías del conjunto de prueba (15 consultas en inglés cada una)}
    \label{tab:met-categorias}
    \small
    \begin{tabular}{|p{3.4cm}|p{10.8cm}|}
        \hline
        \textbf{Categoría} & \textbf{Descripción} \\
        \hline
        Exactas & Sus términos aparecen literalmente en las ofertas relevantes, por ejemplo \textit{\enquote{senior react developer}}. Es el caso favorable a la búsqueda léxica. \\
        \hline
        Desajuste de vocabulario & Describen la función del cargo en lugar de nombrarlo, por ejemplo \textit{\enquote{protect company systems from hackers}} para un cargo de ciberseguridad. Es el caso que motiva este trabajo (Sección~\ref{sec:problema}). \\
        \hline
        Con restricciones & Combinan un perfil con condiciones de salario, modalidad o tipo de contrato. \\
        \hline
        Ruidosas & Tienen errores ortográficos, abreviaturas o una formulación vaga. \\
        \hline
        De agregación & Preguntas estadísticas sobre el catálogo, por ejemplo cuántas ofertas remotas de Python existen. Se evalúan aparte. \\
        \hline
    \end{tabular}
\end{table}

Las decisiones de configuración (modelo de reordenamiento, número de candidatas y umbral de la compuerta de sin resultados) se toman con un conjunto de desarrollo de 20 necesidades, formuladas en 28 consultas (20 en inglés y 8 en español) sin solapamiento con las de prueba, más 10 consultas fuera de dominio que solo sirven para calibrar el umbral; sus resultados no se reportan como resultados del sistema. Para probar la compuerta con consultas que no redactó el autor, una persona sin formación técnica escribe 15 consultas fuera de dominio a partir de una instrucción breve; se congelan antes de correr el pipeline y se verifica en el corpus, sin mirar puntajes, que ninguna oferta las responde.

Se evalúa la configuración de referencia de JobAI, definida en el repositorio e identificada por el \textit{commit} que acompaña a cada corrida; la Tabla~\ref{tab:art-config} la compara con la desplegada y con la evaluada. Durante la evaluación la temperatura de generación es 0 y las respuestas y clasificaciones se guardan en caché, de modo que todos los sistemas usan las mismas clasificaciones. Las consultas se envían sin usuario autenticado, para que la personalización por currículo no se confunda con la recuperación, y las métricas de recuperación se calculan sin la compuerta, cuyo efecto se reporta aparte.

\section{Sistemas comparados}
\label{sec:met-sistemas}

Los trece sistemas de la Tabla~\ref{tab:met-sistemas} producen una lista ordenada de 10 ofertas por consulta sobre el mismo corpus. Los densos recuperan los 50 fragmentos más similares a la consulta y los agregan por oferta con la máxima similitud entre sus fragmentos, la misma regla del pipeline, de modo que la unidad evaluada es la oferta, que es lo que el usuario ve. La comparación entre peldaños consecutivos de la escalera mide el aporte de cada componente, y la del último peldaño con la recuperación densa, el de la Iteración~2 completa. BM25 es la comparación exigente, porque en la evaluación de \citet{Thakur2021} sobre 18 colecciones resultó una referencia que varios modelos densos no superaron; como su procesamiento es propio del inglés, en las consultas en español se lee como una búsqueda léxica sin traducción.

\begin{table}[!htb]
    \centering
    \caption{Sistemas comparados}
    \label{tab:met-sistemas}
    \small
    \begin{tabular}{|p{4.3cm}|p{9.9cm}|}
        \hline
        \textbf{Sistema} & \textbf{Qué representa o aísla} \\
        \hline
        \multicolumn{2}{|l|}{\textit{Líneas base léxicas}} \\
        \hline
        ILIKE & La búsqueda de la plataforma: subcadena, sin distinguir mayúsculas, sobre el título y la empresa \\
        \hline
        BM25 & La línea base léxica fuerte \citep{Robertson2009}; la misma implementación de la búsqueda híbrida (Anexo~\ref{anx:bm25}) \\
        \hline
        \multicolumn{2}{|l|}{\textit{Recuperación densa: sin filtros, sin BM25 y sin reordenamiento}} \\
        \hline
        Denso, pgvector exacto & La referencia para medir el efecto de la búsqueda aproximada \\
        \hline
        Denso, pgvector HNSW & La búsqueda aproximada en pgvector \\
        \hline
        Denso, Qdrant HNSW & El motor, con parámetros de índice idénticos; primer peldaño de la escalera \\
        \hline
        Denso, e5-base-v2 & El modelo de \textit{embeddings} anterior; aproxima la recuperación de la Iteración~1 \\
        \hline
        \multicolumn{2}{|l|}{\textit{Escalera de componentes: cada peldaño agrega un componente al anterior}} \\
        \hline
        + filtros & Las restricciones que el clasificador extrae de la consulta, aplicadas como filtro \\
        \hline
        + híbrida & BM25 bajo el mismo filtro, fusionado con RRF \\
        \hline
        + reordenamiento & El reordenamiento de las 20 primeras candidatas: la configuración completa \\
        \hline
        \multicolumn{2}{|l|}{\textit{Sistemas complementarios}} \\
        \hline
        BM25 + filtros & BM25 con los filtros que el clasificador extrae: la línea base léxica con la misma información que el pipeline \\
        \hline
        BM25 + filtros + reordenamiento & La misma primera etapa léxica con el reordenador de la configuración completa \\
        \hline
        Completa sin búsqueda híbrida & La configuración completa sin BM25: mide la búsqueda híbrida quitándola \\
        \hline
        Completa sin filtros & La configuración completa sin los filtros: mide los filtros quitándolos \\
        \hline
    \end{tabular}
\end{table}

\section{Juicios de relevancia}
\label{sec:met-juicios}

La evaluación sigue el paradigma de Cranfield: un corpus fijo, un conjunto de necesidades de información y juicios de relevancia permiten comparar sistemas de forma repetible y sin usuarios \citep{Sanderson2010}. Los juicios se construyen por agrupamiento (\textit{pooling}) de profundidad 10: para cada consulta se juzga la unión de las 10 primeras ofertas de todos los sistemas, de modo que esas posiciones quedan juzgadas por completo, y las ofertas que ningún sistema recupera en ellas se consideran no relevantes, una aproximación que afecta sobre todo al \textit{recall}. Si se agregan sistemas, se juzgan solo sus pares nuevos.

Cada par necesidad-oferta recibe de un modelo de lenguaje una etiqueta en la escala de cuatro niveles de TREC: 0, no relevante; 1, relacionada; 2, altamente relevante, y 3, perfectamente relevante. El prompt sigue el diseño de UMBRELA \citep{Upadhyay2024}, una reproducción abierta del procedimiento con el que \citet{Thomas2024} mostraron que un LLM puede predecir las preferencias de evaluadores humanos. El juez recibe la necesidad y una ficha estructurada de la oferta, y aplica una regla de \emph{condiciones duras} adaptada al dominio: si la necesidad fija un salario, la modalidad remota, un tipo de contrato o una ubicación, una oferta que lo contradice o no lo informa recibe como máximo la etiqueta 1, para que la métrica distinga a un sistema que respeta las restricciones de uno que las ignora. Como un juez LLM puede tener sesgos sistemáticos, ser sensible al prompt y caer en circularidad si comparte modelo con el sistema evaluado \citep{Faggioli2023,Zheng2023}, el juez es un modelo distinto del generador, un modelo local pequeño solo se usa durante el desarrollo, y el juez final se elige entre tres modelos comerciales con una regla fijada antes de anotar la muestra humana (Anexo~\ref{anx:juez-final}).

Los juicios del juez final se validan contra juicios humanos. De los pares de la corrida final se extrae una muestra aleatoria del 20\,\%, estratificada por categoría de consulta y por la etiqueta del juez disponible al extraerla, con semilla fija, que el autor etiqueta a ciegas con la misma escala y la misma regla. Una cuarta parte de la muestra sirve para calibrar el prompt y elegir el juez final, y las tres cuartas partes restantes se reservan para medir el acuerdo, sin volver a modificar el prompt. El acuerdo se mide con el kappa de \citet{Cohen1960}, sobre la relevancia binaria (etiqueta $\geq 2$) y ponderado cuadráticamente sobre la escala completa, con un intervalo de confianza \textit{bootstrap}, y se considera aceptable un acuerdo al menos sustancial ($\kappa > 0{,}60$) \citep{Landis1977}. Si el juez final no lo alcanza en la relevancia binaria, todos los pares se juzgan de nuevo con un segundo modelo, Claude Sonnet 5.5 (Anthropic), y solo se afirman las conclusiones que se mantienen con ambos. Las etiquetas humanas son de una sola persona; la variación entre evaluadores cambia las cifras absolutas más que la comparación entre sistemas \citep{Voorhees2000}, aunque si un juez LLM puede reemplazar a los humanos sigue en discusión \citep{Soboroff2025,ClarkeDietz2024}.

\section{Métricas}
\label{sec:met-metricas}

La Tabla~\ref{tab:met-metricas} resume las métricas de cada pregunta de evaluación. En las métricas binarias, una oferta es relevante si tiene etiqueta $\geq 2$; las fórmulas de las métricas de PE1 y del kappa están en el Anexo~\ref{anx:formulas}.

\begin{table}[htbp!]
    \centering
    \caption{Métricas por pregunta de evaluación}
    \label{tab:met-metricas}
    \small
    \begin{tabular}{|p{1.9cm}|p{5.3cm}|p{6.9cm}|}
        \hline
        \textbf{Pregunta} & \textbf{Métricas} & \textbf{Procedimiento} \\
        \hline
        PE1 & P@5, P@10, nDCG@10 \citep{Jarvelin2002}, MRR y \textit{recall}@10 (R@10) & Sobre las listas de 10 ofertas, con \texttt{ir\_measures} \citep{MacAvaney2022}; desagregadas por categoría y por idioma \\
        \hline
        PE1, restricciones & Cumplimiento en las posiciones 5 y 10 (C@5, C@10) y exactitud de la extracción & Sin juez: proporción de ofertas devueltas que cumplen los filtros esperados, y de consultas cuyos filtros se extraen exactamente \\
        \hline
        PE1, agregación & Exactitud del enrutamiento y de la respuesta & Frente a una consulta SQL de referencia escrita a mano (Anexo~\ref{anx:agregacion}) \\
        \hline
        PE2 & Fidelidad y relevancia de la respuesta \citep{Es2024} & Implementación de referencia de RAGAS sobre las respuestas de la ruta real del pipeline (Anexo~\ref{anx:ragas}) \\
        \hline
        PE3 & Latencia por etapa, mediana (p50) y percentil 95 (p95) & Etapas locales: 30 repeticiones por consulta y por motor, tras 5 de calentamiento. Generación: 3 repeticiones sobre 30 consultas \\
        \hline
        PE4 & \textit{Recall}@10, latencia, consultas por segundo, completitud y recursos & Ver Sección~\ref{sec:met-benchmark} \\
        \hline
    \end{tabular}
\end{table}

La latencia se mide en cada etapa, desde la clasificación hasta el tiempo hasta el primer \textit{token} generado (TTFT) y el tiempo total; la percibida hasta el primer \textit{token} es la suma, en cada solicitud, de la clasificación, la construcción del contexto y el TTFT. La de la generación, que depende de un servicio externo, se reporta aparte y se mide con los reintentos automáticos del cliente desactivados, para que la espera por límite de uso no cuente como latencia. El retraso que agregan los componentes se contrasta con los umbrales de \citet{Arapakis2014} y \citet{Bai2017}, según los cuales, en la búsqueda web, un retraso añadido menor que 500~ms no se distingue y uno mayor que 1.000~ms se nota con mucha probabilidad. Como se midieron sobre la página de resultados de un buscador y no sobre una respuesta que aparece de forma incremental, son una referencia y no un presupuesto exacto.

\section{Benchmark de almacenes vectoriales}
\label{sec:met-benchmark}

La comparación entre pgvector y Qdrant adapta la metodología de ANN-Benchmarks \citep{Aumuller2020}: para cada valor de $ef_{search} \in \{16, 32, 64, 128, 256, 512\}$ se mide el \textit{recall}@10 frente a la búsqueda exacta y la latencia, y el resultado se presenta como curvas de compromiso en lugar de un único punto de operación. Como el corpus real tiene solo unos miles de fragmentos, se generan además conjuntos sintéticos de 100.000 y 500.000 vectores derivados de los reales, que heredan los atributos de su oferta de origen para que los filtros conserven su correlación con el contenido (Anexo~\ref{anx:sinteticos}). Cada configuración se mide sin filtro y con filtros estructurados de selectividad 1\,\%, 10\,\% y 50\,\%, la búsqueda mixta vector-relacional que estudian \citet{Sanca2024}, y en pgvector con y sin los barridos iterativos del índice. Ambos motores usan el mismo índice HNSW ($m=16$, $ef_{construction}=128$) y se ejecutan uno a la vez con los mismos límites de CPU y memoria; cada medición se repite tres veces y se reporta la mediana, con los vectores de ambos modelos de \textit{embeddings} como prueba de robustez.

\section{Análisis estadístico}
\label{sec:met-estadistica}

La unidad de análisis es la necesidad de información, porque las consultas en español comparten los juicios de su contraparte en inglés y tratarlas como independientes contaría dos veces la misma necesidad. En el análisis global y por categoría, la puntuación de cada una de las 60 necesidades que producen una lista de ofertas es la media de sus formulaciones; los análisis por idioma usan las consultas de ese idioma. Cada diferencia se contrasta con la prueba de aleatorización pareada (10.000 permutaciones), que \citet{Smucker2007} recomiendan frente a la de signos y la de Wilcoxon, y se acompaña de un intervalo de confianza \textit{bootstrap} del 95\,\%. nDCG@10 es la métrica principal, porque usa las etiquetas graduadas y penaliza el orden; sus valores $p$ se ajustan con el procedimiento de \citet{Holm1979} dentro de dos familias fijadas antes de la corrida final: la de cada sistema frente a cada línea base (14 comparaciones) y la de los peldaños de la escalera (9), que compara el modelo multilingüe con e5-base-v2 (en todas las consultas, en inglés y en español), pgvector con HNSW con su búsqueda exacta, Qdrant con pgvector, cada componente con el peldaño anterior y la configuración completa con la recuperación densa. Las comparaciones de los sistemas complementarios forman una tercera familia, de cinco. El análisis es comparativo y descriptivo, sin hipótesis a priori; las demás métricas se reportan con su intervalo y su valor $p$ sin ajustar, y el nivel de significancia es $\alpha = 0{,}05$.

\section{Amenazas a la validez y repetibilidad}
\label{sec:met-validez}
\label{sec:met-reproducibilidad}

La Tabla~\ref{tab:met-validez} resume las amenazas principales y cómo se mitigan.

\begin{table}[!htb]
    \centering
    \caption{Amenazas a la validez y su mitigación}
    \label{tab:met-validez}
    \footnotesize
    \begin{tabular}{|p{2cm}|p{6.3cm}|p{5.9cm}|}
        \hline
        \textbf{Tipo} & \textbf{Amenaza} & \textbf{Mitigación} \\
        \hline
        Interna & Sesgos del juez automático, incluido el que favorece a los sistemas que ordenan con LLM \citep{Frobe2025}. & Validación humana, kappa y plan B con un segundo juez; juez final de otra familia que el generador y el reordenador, lo que no elimina el sesgo. \\
        \hline
        Interna & El juez local es el modelo que analizó el corpus y redactó el resumen que indexa la búsqueda densa. & Solo se usa en el desarrollo; el juez final es de otra familia. \\
        \hline
        Interna & El agrupamiento favorece a los sistemas que lo originan. & Todos los sistemas aportan al agrupamiento; si se agregan sistemas, se juzgan sus pares nuevos y se recalculan las cifras. \\
        \hline
        Interna & Un solo anotador humano y consultas redactadas por el autor. & Validación a ciegas después de elegir el juez; categorías fijadas antes; las consultas fuera de dominio de prueba las escribe otra persona. \\
        \hline
        Interna & Decisiones tomadas después de ver resultados: la Iteración~2, los criterios de evaluación y los análisis complementarios. & Configuración solo con el conjunto de desarrollo; familia de pruebas aparte para los complementarios, que no cambian el sistema; las fallas vistas en la prueba no se corrigen. \\
        \hline
        Externa & Un dominio, un idioma de las ofertas, dos modelos de la misma familia y un corpus real pequeño. & No se generaliza; el \textit{benchmark} sintético complementa las conclusiones sobre escala. \\
        \hline
        Medición & La latencia incluye el protocolo y la biblioteca cliente de cada motor, y la del reordenamiento depende de la GPU. & Se interpreta junto con la escala de los datos y el equipo de evaluación. \\
        \hline
        Constructo & La ganancia de los filtros en nDCG@10 depende en parte de la regla de condiciones duras, y la relevancia juzgada no equivale a la utilidad real. & El cumplimiento de restricciones (C@5) se calcula sin juez; la evaluación con usuarios queda fuera del alcance (Sección~\ref{sec:alcance}). \\
        \hline
        Confiabilidad & Los LLM no son deterministas y pueden retirarse sin aviso. & Cada respuesta y veredicto se guarda con el modelo exacto y la fecha. \\
        \hline
    \end{tabular}
\end{table}

Cada corrida se ejecuta con la interfaz de línea de comandos del sistema y registra el \textit{commit}, los modelos y parámetros, el resumen SHA-256 del corpus y los resultados por consulta. Las versiones del software y de los modelos están fijadas (pgvector 0.8.2, Qdrant 1.19.1), y los veredictos, respuestas y clasificaciones se guardan en caché, de modo que el análisis se repite sin nuevas llamadas a modelos de lenguaje y el mismo análisis sobre los mismos datos da las mismas cifras (\emph{repetibilidad}). El código, las consultas, los juicios y los resultados de cada corrida se publican en el repositorio del trabajo (Anexo~\ref{anx:repositorio}); el corpus no se publica, por los términos de uso de las fuentes, de modo que la \emph{reproducibilidad} por terceros es parcial: pueden repetir el análisis con los juicios y las listas publicadas y reconstruir un corpus con el procedimiento documentado, pero no el mismo corpus.
